% !TEX root = mockup.tex
% =========================================================================
% MOCKUP CHAPTER 0 -- Preface
% Purpose: verify that an UNDECORATED chapter (no cover page, no skills box,
% no take-home box, unnumbered) still gets the correct body layout, headers,
% and float behaviour. Blind text throughout.
% =========================================================================
\startchapterstar{Prologue:\\From Raindrops to River Systems}
\chaptershortname{Prologue}

% =========================================================================
% FIRST PAGE -- title row + skills box at the top of column 2
% =========================================================================
\chapterfirstpage{}

% --- p1: plain opening, no decoration ------------------------------------
\section{A prologue has no cover page}

\lipsum[1]

Note that this chapter opens directly on a text page: no full-bleed cover,
no chapter-number bar, no professional-skills box. The running header still
works, and the page number still sits on the outer edge.

% --- p2: S figure top + S figure bottom ----------------------------------
\figplaceholder[L][b!]{C00-mtx-processes}{C00-mtx-processes}{Two-column figure, bottom of
page. Image block is 172\,mm wide, \lipsum[2]}

\lipsum[6-8]

% --- p3: L figure top -----------------------------------------------------
\section{Two-column figures in an undecorated chapter}

\figplaceholder[L][t]{C00-mtx-cycles}{C00-mtx-cycles}{Two-column figure, top of
page. Image block is 172\,mm wide. \lipsum[2]}

\lipsum[8-12]

\figplaceholder[S][t]{C00-mtx-interactions}{C00-mtx-interactions}{One-column figure, top of
page. Image block is 83.5\,mm wide. \lipsum[2]}


% --- p4: M figure with side caption --------------------------------------
\figplaceholder[M][t]{C00-mtx-ranges}{C00-mtx-ranges}{Medium figure with the
caption set beside it on the right. Image block is 126\,mm by 100\,mm, and
the caption takes the remaining 41\,mm of the text width. This caption is
deliberately long so the side column wraps to several lines and you can see
how tall a side caption is allowed to get before it looks wrong.}

\lipsum[11-15]

% --- p5: table (always two-column, top) ----------------------------------
\section{Tables}

\lipsum[11-15]

\begin{table*}[t!]
\caption{Tables are always two-column and always placed at the top of a page.
This one is deliberately wide enough to test the full text measure.}
\label{tab:mock-preface}
\sffamily\scriptsize
\begin{tabular}{p{0.16\textwidth} p{0.16\textwidth} p{0.20\textwidth} p{0.20\textwidth} p{0.18\textwidth}}
\toprule
\textbf{Column one} & \textbf{Column two} & \textbf{Column three} &
\textbf{Column four} & \textbf{Column five} \\
\midrule
Lorem ipsum dolor & Sit amet consectetur & Adipiscing elit sed do eiusmod &
Tempor incididunt ut labore et dolore & Magna aliqua ut enim \\
\addlinespace
Ad minim veniam & Quis nostrud & Exercitation ullamco laboris nisi &
Ut aliquip ex ea commodo consequat & Duis aute irure dolor \\
\addlinespace
In reprehenderit & In voluptate velit & Esse cillum dolore eu fugiat &
Nulla pariatur excepteur sint occaecat & Cupidatat non proident \\
\addlinespace
Sunt in culpa & Qui officia deserunt & Mollit anim id est laborum &
Sed ut perspiciatis unde omnis iste & Natus error sit voluptatem \\
\bottomrule
\end{tabular}
\end{table*}

\lipsum[14-16]

% --- p6: equations, cross-references, lists ------------------------------
\section{Equations, cross-references and lists}

\lipsum[17]

A numbered equation, which should carry the chapter prefix:

\begin{equation}
\bar{Q} = \bar{P} - \bar{E} - \frac{\mathrm{d}S}{\mathrm{d}t}
\label{eq:mock-balance}
\end{equation}

Cross-reference tests, all of which must resolve:
\Cref{eq:mock-balance}, \Cref{C00-mtx-processes}, \Cref{C00-mtx-ranges} and
\Cref{tab:mock-preface}.

\begin{itemize}
\item An itemised list inside the body text, first item.
\item A second item, slightly longer so that it wraps onto a second line in
      an 83.5\,mm column and you can inspect the hanging indent.
\item A third item.
\end{itemize}

\begin{enumerate}[(a)]
\item An enumerated list using the \texttt{shortlabels} option.
\item Second entry.
\end{enumerate}

\lipsum[18]

\subsection{A subsection heading}

\lipsum[19]

\subsubsection{A subsubsection heading}

\lipsum[20]

\paragraph{A run-in paragraph heading} \lipsum[21]

An author note for editorial review: \authornote{this is what an author note
looks like inline in the body text.}

\lipsum[22]

% --- p7: no take-home box in a preface -----------------------------------
\section{The preface ends without a take-home box}

\lipsum[23-24]